\documentclass[10pt]{article}
\usepackage[utf8]{inputenc}
\usepackage[T1]{fontenc}
\usepackage{amsmath}
\usepackage{amsfonts}
\usepackage{amssymb}
\usepackage[version=4]{mhchem}
\usepackage{stmaryrd}
\usepackage{bbold}
\usepackage{graphicx}
\usepackage[export]{adjustbox}
\graphicspath{ {./images/} }
\usepackage{mathrsfs}

\begin{document}
Установлено, что О. п. Ф: $S \rightarrow M_{G}$ горизонтально (см. [1], [3]). Особенности О. п. описываются т е о р ем о й Ш ми да о и ль п о тентной о р б и те, к-рая в случае, когда $S=\bar{S} \backslash\{0\}-$ кривая с выколотой точкой, утверждает, что если $z-$ локальная координата на $S, z(0)=0$, то при $z \rightarrow 0 \Phi(z)$ асимптотически близко ік

\[
\exp \left(\frac{\log z}{2 \pi i} N\right) a,
\]

где $a \in \breve{X}_{G}$, а $N \in\left(\mathrm{Si}_{\mathrm{Q}}\right.$ - нильпотөнтный элемент (см. [4]). Относительно группы монодромии

\[
\Phi_{*}\left(\pi_{1}(S, s)\right) \subset G_{\mathbb{Z}}
\]

известно, что ее образ полупрост во всяком рациональном представлении группы $G$, а преобразования обхода $T$ вокруг дивизора с нормальными пересечениями $\bar{S} \backslash S$ в гладкой компактификкации $\bar{S}$ многообразия $S$ порождают к в а 3 и у ни по т е н т ны е (т. е. имеющие в качестве собственных значений корни из 1) әлементы $\Phi_{*}(T) \in G_{\mathbb{Z}} \cdot$ Важность группы монодромии подчеркивает те о р е м ж ж с т к о т и (см. [1], [2], [4]): если над $S$ имеются два семейства алгебраич. мноғообразий, то соответствующие О. п. $\Phi_{1}$ и $\Phi_{2}$ из $S$ в $M_{G}$ совпадают тогда и только тогда, когда $\Phi_{1}\left(s_{0}\right)=$ $=\Phi_{2}\left(s_{0}\right)$ в нек-рой точке $s_{0} \in S$ и гомоморфизмы $\Phi_{i *}:$ $\pi_{1}\left(S, s_{0}\right) \rightarrow G_{\mathbb{Z}}, i=1,2, \quad$ совпадают.

Законченные результаты о строении ядра и образа О. П. относятся в основном к случаям кривых и КЗповерхностей. Если $\left\{X_{s}\right\}$ - семейство многообразий указанного типа и $\Phi(s)=\Phi\left(s^{\prime}\right)$, то $X_{s \rightarrow}^{\sim} X_{s^{\prime}}$ (т е о р е м а Т о р е л л), а для $K 3$-поверхностей максимально возможный образ О. п. совпадает с $M_{G}$ (см. [7]). В случае кривых образ О. п. частично описан (соотношения Шоттки - Юнга, см. [6], [8]). Существует г и п о т е з а $\Gamma$ р и фф и т с а о том, что многообразие модулей допускает частичную аналитич. компактификацию, т. е. открытое вложение в такое аналитич. пространство $\bar{M}_{G}$, что O. п. $S \rightarrow M_{G}$ продолжается до голоморфного отображения $\bar{S} \supset S$ для всякой гладкой компактификации $S \supset S$. Такая компактификация известна (1983) лишь для случая, когда $X_{G}$ - симметрическая область $[9]$

Лum.: [1] r p и фо фи и с $\Phi$. А., "Успехи матем. наук", 1970, T. 25 , B. 3, c. $175-234 ;$ [2] G r i f f i t h s Ph. A., в \%н.: Actes du 1971 , p. $113-19 ;$ [3] D e 1 i g n e P., Travaux de Griffiths, в кH.: 1971 , p. 113-19; [3] D e 1 i g n e P., Travaux de Griffiths, в кн.:
Seminaire Bourbaki. $1969 / 70$, B.- N. Y. Hdlb., 1971, p. $213-35$ [4] S c h m i d W., "Invent. math.", 1973, v. 22, p. 211-319; [5] C a t t a n i E. H., K a p l a n A. G., "Duke Math. J.", 1977,
v. $44, N_{2} 1$, p. 1-43; [6] Д у б р о в и н Б. А., "Успехи матем. наук", 1981 , т. 36, в. 2 , с. $11-80 ;$ [7] К у л и к о в В. А., там же, 1977 , т. 32 , в. 4, с. $257-58 ;$ [8] М а м ф о р д Д., “Математика", 1973 , т. $17, \mathcal{N}_{2}$, c. $34-42 ;$ [9] B a i l y W., B o r e 1 A., "Ann. Math.", $1966, \mathrm{v} .84, \mathrm{p} .442-528$. И. Овсеевич.
ОТОБРАЖЕНЙ КЛАССБІ - важнейшие классы негрерывных отображений, рассматриваемые в обшей топологии и ее приложениях. $\mathrm{K}$ ним относятся: о тк ры ты е о тоб р а же и я - такие, что образ любого открытого множества является открытым множеством; 3 а м к н у т ы е о т о б р а же ни я- такие, при к-рых образ каждого замкнутого множества замкнут; б и о м п а к т н ы е о тоб р а ж н н я для них прообраз любой точки является бикомпактным множеством; с о в е р II $\boldsymbol{\theta}$ н ны е о т о б р аж е н и я - замкнутые бикомпактные отображения. Ф а т т р ны е о т о б р а же н и я определяются требованием: множество в образе открыто в том и только в том случае, если его полный прообраз открыт. Важны также открытые бикомпактные отображения,

\includegraphics[max width=\textwidth, center]{2023_07_08_8d638532475d349ab56dg-1}
отоб ражев в у п л о т н е в и я - последние определяютея как взаимно однозначные непрерывные отображения на. Таким образом, при классификации отображениї̈ в общей топологии ограничөния накладываются либо на поведөние (при перөходе к образу) открытых или замкнутых множеств, либо на свойства прообразов множеств. Второй подход приводит, в частности, к следующим $О$. к. М о но то н ны е о то б р а ж ен и я-те, при к-рых прообраз каждой точки нульмерен. Конечнок р а т ны е отоб р а е ния характеризуются конөчностью всех прообразов точек. Отображения, при к-рых прообраз каждого бикомпактного множества бикомпактен, наз. $k-0$ т о б р aж е в и я м. Соединением ограничөний первого и второго типа выделяются основныө классы непрерывных отображений в общей топологии. Самими определениями $\mathrm{O} . \kappa$. естественно организуются в нек-рую иерархию, к-рая может быть положена в основу систематич. классификации топологич. пространств [1]. Эта классификация строится как результат решения вопросов следующих двух типов. Дан класс пространств $\mathcal{A}$, целесообразность выделения к-рого не вызывает сомнений, и пусть $\mathscr{B}-$ нек-рый класс отображений из нашей исходной иерархии. Требуется охарактеризовать посредством внутренних топологит. инвариантов образы пространств из класса $\mathcal{A}$ при всевозможных отображениях из класса $\mathscr{B}$. Вопросы второго типа аналогичны - требуется охарактеризовать прообразы пространств из класса $\mathcal{A}$ при отображениях из класса $\mathscr{B}$.

При решении вопросов указанных двух типов получаются совсем не очевидные теоремы общего характера. Напр., пространства с первой аксиомой счетности - это в точности образы метрич. пространств при непрерывных открытых отображениях. Пространства с равномерной базої и только они - образы метрич. пространств при открытых бикомпактных отображениях. Пространства Фреше - Урысона характеризуются как псевдооткрытые образы метрич. пространств, а секвенциальные пространства - это факторпространства метрич. пространств. Далее, прообразы метрич. пространств при совершенных отображения - это в точности паракомпактные перистые пространства, а прообразами полных метрич. пространств являются паракомпактные пространства, полные по Чеху. Непрерывные образы пространств со счетной базой - пространства со счетной сетью. При систематич. следовании указанному пути получается единая взаимная классификация пространств и отображений.

Особую роль среди различных классов непрерывных отображений занимает класс факторных отображений. Важнейпей особенностью факторных отображений является то, что они могут служить средством для построения новых топологич. пространств. А именно, если дано отображение $f$ топологич. пространства $X$ на нек-рое множество $Y$ (напр., если рассматривается естественное отображение $\pi$ пространства $X$ на множество всех элементов нек-рого разбиения этого пространства), то на множестве $Y$ всегда можно ввести естественную топологию требованием, чтобы отображение $f$ было факторным: множество $V \subset Y$ объявляется открытым в том и только в том случае, если его полный прообраз $f^{-1}(V)$ открыт в пространстве $X$. Помимо уже названных, весьма важны неприводимье отображения, напр. в теории абсоліотов. См. также Бибакторное отображение, Многозначное отображение.

Лum.: [1] А p х а н г е ль с к и й А. В., “Успехи матем. наук», 1966, т. 21, в. 4, с. 133-84. A. В. Архангельский.

ОТОБРАЖЕНИИ МЕТОД - то Же, что изображений метод.


\end{document}